\writeSectionFrame{\presentationTitle}{Erbauer}
\begin{frame}{Steckbrief}
	\begin{itemize}
 		\item Deutscher Name: Erbauer
 		\item Englischer Name: Builder
 		\item Gruppe: Erzeugungsmuster
 	\end{itemize}
\end{frame}

\begin{frame}{Erbauer}
	\begin{block}{Zweck}
 		Trenne die Konstruktion eines komplexen Objekts von seiner Repräsentation, so dass deselbe Konstruktionsprozess unterschiedliche Repräsentationen erzeugen kann.
 	\end{block}
\end{frame}

\begin{frame}{Allgemeines Klassendiagramm}
    \begin{tikzpicture}[scale=0.7, transform shape]
        \umlclass[x=17,y=-3]{Product}{
            +attribute1: Type\\+attribute2: Type}{}
        
        % Klasse Builder
        \umlclass[x=12,y=0]{Builder}{}{
            +setAttribute1(attr1: Type) : Builder\\
            +setAttribute2(attr2: Type) : Builder\\+build() : Product}{}

        \umlclass[x=12,y=-3]{ConcreteBuilder}{}{}
        
        % Klasse Director
        \umlclass[x=5,y=0]{Director}{
            +construct(builder: Builder)\\
            -builder: Builder}{}
    
        % Beziehungen
        \umlinherit{ConcreteBuilder}{Builder}
        \umlaggreg{Director}{Builder}
        \umluniassoc{ConcreteBuilder}{Product}
    \end{tikzpicture}
\end{frame}

\frameWithImageOnly{\BILDERPFAD/erbauer1.jpg}

\begin{frame}{Klassendiagramm Beispiel Häuserbau}
    \begin{tikzpicture}[scale=0.7, transform shape]
        \umlclass[x=17,y=-4]{House}{
        }{}
        
        % Klasse Builder
        \umlclass[x=12,y=0]{Builder}{}{
            +addSwimmingPool()\\
            +addGarden() \\
            +addRoom()\\
            +setWallThickness(thickness: int)\\
            +buildHouse(): House
            }{}

        \umlclass[x=12,y=-4]{HouseBuilder}{}{}
        
        % Klasse Director
        \umlclass[x=5,y=0]{Director}{
        }{}
    
        % Beziehungen
        \umlinherit{HouseBuilder}{Builder}
        \umluniassoc{Director}{Builder}
        \umluniassoc{HouseBuilder}{House}
    \end{tikzpicture}
\end{frame}

\begin{frame}[fragile]{Erbauer}
	\begin{exampleblock}{Builder.java}
 		\begin{lstlisting}
public interface Builder {
	void addSwimmingPool();
	void addGarden();
	void addGarage();
	void addRoom();
	void setWallThickness(int wallThickness);
	void setHouseType(HouseType houseType);
	House buildHouse();
	void newHouse();
}
 		\end{lstlisting}
 	\end{exampleblock}
\end{frame}

\begin{frame}[fragile]{Konkreter Erbauer}
	\begin{exampleblock}{HouseBuilder.java}
        \begin{lstlisting}
public class HouseBuilder implements Builder {
	private int wallThickness;
	private boolean swimmingPool;
	private boolean garden;
	private int garageCount;
	private int roomCount;
	private HouseType houseType;
	
	@Override
	public void newHouse() {
		wallThickness = 0;
		swimmingPool = false;
		garden = false;
		garageCount = 0;
		roomCount = 0;
	}
        \end{lstlisting}
 	\end{exampleblock}
\end{frame}

\begin{frame}[fragile]{Konkreter Erbauer}
	\begin{exampleblock}{HouseBuilder.java}
 		\begin{lstlisting}
@Override
public void addSwimmingPool() {
    swimmingPool = true;
}

@Override
public void addGarden() {
    garden = true;
}

@Override
public void addGarage() {
    garageCount++;
}
 		\end{lstlisting}
 	\end{exampleblock}
\end{frame}

\begin{frame}[fragile]{Direktor}
	\begin{exampleblock}{HouseDirector.java}
 		\begin{lstlisting}
public class HouseDirector {
	public void constructSingleFamilyHouse(Builder builder) {
		builder.newHouse();
		builder.addGarage();
		for (int i = 0; i < 4; i++) {
			builder.addRoom();
		}
		
		builder.addSwimmingPool();
		builder.addGarden();
		builder.setHouseType(HouseType.SINGLE_FAMILY_HOUSE);
		builder.setWallThickness(2);
	} 		    
 		\end{lstlisting}
 	\end{exampleblock}
\end{frame}

\frameWithImageOnly{\BILDERPFAD/erbauer2.jpg}
\note{
    Im nächsten Beispiel simulieren wir einen Programm, mit dem der Benutzer sich einen PC zusammenstellen und dann bestellen kann. 
}

\begin{frame}{Klassendiagramm Beispiel Computerbau}
    \begin{tikzpicture}[scale=0.7, transform shape]
        \umlclass[x=6,y=-4]{Computer}{
        }{}
        
        % Klasse Builder
        \umlclass[x=12,y=0]{ComputerBuilder}{}{
	       +setProcessor(processor: String)\\
          +setRam(ram: int)\\
	       +setHardDrive(hardDrive: int)\\
          +setGraphicsCard(graphicsCard: String )\\
            }{}

        \umlclass[x=12,y=-4]{ComputerBuilderImpl}{}{}
        
        % Klasse Director
        \umlclass[x=3,y=0]{ComputerDirector}{
        }{}
    
        % Beziehungen
        \umlinherit{ComputerBuilderImpl}{ComputerBuilder}
        \umluniassoc{ComputerDirector}{ComputerBuilder}
        \umluniassoc{HouseBuilder}{Computer}
    \end{tikzpicture}
\end{frame}

\begin{frame}[fragile]{Erbauer}
	\begin{exampleblock}{Builder.java}
 		\begin{lstlisting}
public interface ComputerBuilder {
	void setProcessor(String processor);
	void setRam(int ram);
	void setHardDrive(int hardDrive);
	void setGraphicsCard(String graphicsCard);
	Computer build();
}
 		\end{lstlisting}
 	\end{exampleblock}
\end{frame}

\begin{frame}[fragile]{Direktor}
	\begin{exampleblock}{ComputerDirector.java}
 		\begin{lstlisting}
public class ComputerDirector {
    private ComputerBuilder builder;

    public Computer constructGamingComputer() {
    	builder.setProcessor("Intel Core i9");
    	builder.setRam(32);
    	builder.setHardDrive(1024);
    	builder.setGraphicsCard("NVIDIA GeForce RTX 4090");
        return builder.build();
    }

    public Computer constructOfficeComputer() {
    	builder.setProcessor("Intel Core i5");
    	builder.setRam(8);
    	builder.setHardDrive(256);
        builder.setGraphicsCard("Integrated");
        return builder.build();
    }
}    
 		\end{lstlisting}
 	\end{exampleblock}
\end{frame}

\frameWithImageOnly{\BILDERPFAD/erbauer3.jpg}
\note{
	Wir erweitern unseren Häusererbauer um ein weiteres Erbauermuster - und zwar sollen die	Daten der Häuser in Dateien gespeichert werden: in csv- und xml-Dateien.
}

\begin{frame}{Klassendiagramm Beispiel Häuserbau mit Dateien}
    \begin{tikzpicture}[scale=0.7, transform shape]
        \umlclass[x=17,y=-4]{File}{
        }{}

        \umlclass[x=14,y=-6]{XmlFile}{
        }{}

        \umlclass[x=14,y=-2]{CsvFile}{
        }{}
        
        % Klasse Builder
        \umlclass[x=10,y=0]{FileBuilder}{}{
        }{}

        \umlclass[x=10,y=-4]{XmlFileBuilder}{}{}

        \umlclass[x=14,y=-4]{CsvFileBuilder}{}{}
        
        % Klasse Director
        \umlclass[x=5,y=0]{FileDirector}{
        }{}
    
        % Beziehungen
        \umlinherit{XmlFileBuilder}{FileBuilder}
        \umlinherit{CsvFileBuilder}{FileBuilder}

        \umlinherit{XmlFile}{File}
        \umlinherit{CsvFile}{File}
        
        \umluniassoc{FileDirector}{FileBuilder}
        \umluniassoc{CsvFileBuilder}{CsvFile}
        \umluniassoc{XmlFileBuilder}{XmlFile}
    \end{tikzpicture}
\end{frame}

\begin{frame}[fragile]{Erbauer}
	\begin{exampleblock}{FileBuilder.java}
 		\begin{lstlisting}
public abstract class FileBuilder {
	protected File file;
	
	public abstract void appendData(House house);

	public File getFile() {
		return file;
	}
}
 		\end{lstlisting}
 	\end{exampleblock}
\end{frame}

\begin{frame}[fragile]{Konkreter Erbauer}
	\begin{exampleblock}{XmlFileBuilder.java}
        \begin{lstlisting}
public class XmlFileBuilder extends FileBuilder {
	@Override
	public void appendData(House house) {
		file.addLine("<house houseType=\"" + house.getHouseType() 
            + "\">\n");
		file.addLine("<roomCount>" + house.getRoomCount() 
            + "</roomCount>\n");
		file.addLine("<garageCount>" + house.getGarageCount() 
            + "</garageCount>\n");
		file.addLine("<garden>" + house.hasGarden() + "</garden>\n");
		file.addLine("<swimmingpool>" + house.hasSwimmingPool() 
            + "</swimmingpool>\n");
		file.addLine("<garden>" + house.hasGarden() + "</garden>\n");
		file.addLine("</house>");
	}
}            
        \end{lstlisting}
 	\end{exampleblock}
\end{frame}

\begin{frame}[fragile]{Konkreter Erbauer}
	\begin{exampleblock}{CsvFileBuilder.java}
 		\begin{lstlisting}
public class CsvFileBuilder extends FileBuilder {
	private static final String SEPARATOR = ";";
	
	public CsvFileBuilder() {
		file = new CsvFile();
		file.addLine("");
	}
	
	@Override
	public void appendData(House house) {
		file.addLine(house.getHouseType() + SEPARATOR +
				house.getRoomCount() + SEPARATOR + 
				house.getGarageCount() + SEPARATOR +
				house.getWallThickness() + SEPARATOR +
				house.hasGarden() + SEPARATOR +
				house.hasSwimmingPool() + "\n");
				
	}
}
        \end{lstlisting}
 	\end{exampleblock}
\end{frame}

\begin{frame}[fragile]{Direktor}
	\begin{exampleblock}{FileDirector.java}
 		\begin{lstlisting}
public class FileDirector {
	public void createFile(List<House> houses, 
			FileBuilder fileBuilder) {
		for (House house: houses) {
			fileBuilder.appendData(house);
		}
	}
}
 		\end{lstlisting}
 	\end{exampleblock}
\end{frame}

\begin{frame}[fragile]{Buildermuster 1 zum Bauen von Häusern}
	\begin{exampleblock}{FileWriterClient.java}
 		\begin{lstlisting}
private static List<House> createHouseList() {
    final List<House> houses = new ArrayList<>();
    Director director = new Director();
    HouseBuilder houseBuilder = new HouseBuilder();

    director.constructSkyscraper(houseBuilder);
    House mySkyscraper = houseBuilder.buildHouse();
    houses.add(mySkyscraper);
    
    director.constructAppartmentHouse(houseBuilder);
    House myAppartmentHouse = houseBuilder.buildHouse();
    houses.add(myAppartmentHouse);
    
    director.constructSingleFamilyHouse(houseBuilder);
    House mySingleFamilyHouse = houseBuilder.buildHouse();
    houses.add(mySingleFamilyHouse);
    
    return houses;
} 		    
 		\end{lstlisting}
 	\end{exampleblock}
\end{frame}

\begin{frame}[fragile]{Buildermuster 2 zum Schreiben der Dateien}
	\begin{exampleblock}{FileWriterClient.java}
 		\begin{lstlisting}
public static void main(String[] args) {
    FileDirector fileDirector = new FileDirector();
    FileBuilder xmlFileBuilder = new XmlFileBuilder();
    FileBuilder csvFileBuilder = new CsvFileBuilder();
    
    List<House> houses = createHouseList();
    fileDirector.createFile(houses, xmlFileBuilder);
    fileDirector.createFile(houses, csvFileBuilder);
    System.out.println(csvFileBuilder.getFile().getContent());
    System.out.println(xmlFileBuilder.getFile().getContent());
} 		    
 		\end{lstlisting}
 	\end{exampleblock}
\end{frame}